\bookStart{Speeches of Allwise}[Al·víss mǫ́l]
\setBookCode{Alvissmal}

\begin{flushright}%
\textbf{Dating} \parencite{Sapp2022}: C10th (0.851)

\textbf{Meter:} \Ljodahattr%
\end{flushright}

\section{Introduction}

The \textbf{Speeches of Allwise} (Alvissmal) is a list of poetic synonyms set in a frame narrative of Thunder being visited by a dwarf insisting that he has been promised his daughter’s hand.

\subsection{Textual preservation and transmission}

\Alvissmal\ is found in whole in \Regius\ on foll.~19v/7--20r/19, where it follows \textlink{Volundarkvida} and is followed by \textlink{HelgakvidaOne}.  It concludes the mythological group of poems (to which also \textlink{Volundarkvida} belongs) in the manuscript; the immediately following \textlink{HelgakvidaOne} introduces the heroic group.

The poem was known to Snorre.  Sts.~20 and 30 are cited in \Skaldskaparmal, where they are attributed to \emph{Al·víss mǫ́l} in TODO (A, C) and to \emph{Al·svinns mǫ́l} ‘the Speeches of Allswith’ in all other mss. (TODO).

\subsection{Origin and motifs}

There is not much of a story in \Alvissmal, but two elements of the frame narrative (sts.~1--8, 35) are of interest.

The first is the motif of an enemy of the Gods wishing to steal away one of their women.  This is also attested in the myth of the ettin Master Builder (\textlink{Voluspa}[24] nn., \Gylfaginning~42) and in \textlink{Thrymskvida} (see Introduction to that poem), where the ettin Thrim demands the hand of the goddess Frow as a ransom for Thunder’s hammer Millner.

The second motif is that of a dwarf perishing (usually by turning to stone, e.g.~\textlink{HelgakvidaHjorvardssonar}[29]--30, although that is not stated explicitly in \Alvissmal) when exposed to sunlight at daybreak, a motif famously employed by J.R.R Tolkien in \emph{The Hobbit}.

\subsubsection{The language of the Gods}

The bulk of the poem (sts.~9--34) is however taken up by the listing of poetic synonyms for various phenomena.  \Alvissmal\ contrasts the normal terms used by men with those used by the Elves, Dwarfs, Gods, and those in Hell or the Underworld.
Some of these synonyms, especially those said to be used by the Gods, are quite archaic and represent older common Indo-European and Germanic words that have been displaced by younger words in the common register (10/1 \emph{fold}, 26/1 \emph{funi}, 32/1 \emph{barr}) or alternate forms of the same IE root (16/1 \emph{sunna}).  Others are transparent kennings not found elsewhere.

The fascinating idea of a language of the Gods is not unique to the Norse but is also attested in other cultures.

Notably, Homer often mentions that things have different names among Gods and mortals.
So a certain hill at Troy “do men truly call ‘Thorny Hill’ (\textgreek{Βατίειαν}), but the Undying ‘the Barrow of much-springing Myrina’ (\textgreek{σῆμα πολυ·σκάρθμοιο Μυρίνης})” (\Iliad\ 2.813--814); a certain bird “do Gods call Bronzen (\textgreek{χαλκίδα}), but men \textgreek{κύμινδιν}” (\Iliad\ 14.291); and the river Scamander “do Gods call Xanthus (\textgreek{Ξάνθον}), but men Scamander (\textgreek{Σκάμανδρον})” (\Iliad\ 20.74).

\subsection{Meter and stanza}

The meter of \Alvissmal\ is \Ljodahattr.  The poem consists of 140 lines (71 full-lines, 70 long-lines) divided into 35 sts., all of which follow the typical structure \emph{abcabc} except for the final st.~35 which has the 5-line structure \emph{abcabcc}.

\subsection{Summary}

As mentioned above the frame narrative of the poem consists of sts.~1--8 at the beginning and 35 at the very end.

The section dealing with names (sts.~9--34) is tightly structured and the phenomena are frequently paired according to their nature, e.g. earth and heaven (9--12), sun and moon (15--18), water and fire (23--26).  The following phenomena are dealt with:

\begin{small}
Sts.~9--10: earth, 11--12: heaven, 13--14: moon, 15--16: sun, 17--18: clouds, 19--20: wind, 21--22: calm weather, 23--24: water, 25--26: fire, 27--28: wood, 29--30: night, 31--32: grain, 33--34: ale.
\end{small}

\newpage

\section{Text}

\subsection{The Speeches of Allwise}

\bvg\bva\mssnote{\Regius~19v/7--8}%
\speakernote{[All·víss kvað:]}%
\llap{„}\alst{B}ękki \alst{b}ręiða \hld\ nú skal \alst{b}rúðr með mér &
\ind hęim ï \alst{s}inni \alst{s}núask; &
\alst{h}ratat of mę́gi \hld\ mun \alst{h}vęrjum þikkja; &
\ind \alst{h}ęima skal⸗at \alst{h}víld nema.“\eva

\bvb\speakernoteb{[Quoth Allwise:]}%
“{\huge S}\textsc{pread out on the benches} shall now the bride with me, \\
\ind {[and]} turn home by my side. \\
A hurried engagement it will seem to each; \\
\ind at home shall she not take rest!”\evb\evg


\bvg\bva\mssnote{\Regius~19v/9--10}%
\speakernote{[Þórr kvað:]}%
\llap{„}Hvat ’s þat \alst{f}ira; \hld\ hví est svá \alst{f}ǫlr umb nasar; &
\ind vast-u ï \alst{n}ǫ́tt með \alst{n}á? &
\alst{Þ}ursa líki \hld\ þikki mér ȧ \alst{þ}ér vesa; &
\ind est⸗at-tu til \alst{b}rúðar \alst{b}orinn.“\eva

\bvb\speakernoteb{[Quoth Thunder:]}%
“What sort of man is this; why art thou so pale about the nose; \\
\ind wast thou tonight with a corpse? \\
The likeness of a thurse methinks thou hast; \\
\ind thou wast not born for this bride!”\evb\evg


\bvg\bva\mssnote{\Regius~19v/10--12}%
\speakernote{[All·víss kvað:]}%
\llap{„}\alst{A}l·víss ek hęiti \hld\ bý’k fyr \alst{jǫ}rð neðan &
\ind á’k undir \alst{st}ęini \alst{st}að. &
\edtrans{\alst{v}agna \alst{v}ers}{man of wagons}{\Bfootnote{The “wagons” may here be constellations in the heavens, namely the \emph{Charles’ Wain} (Great Bear, Big Dipper) and \emph{Women’s Wain} (Little Bear, Little Dipper).  Cf.~\Skaldskaparmal\ 31, where heaven/the sky is kenned \emph{land sólar ok tungls ok himin-tungla, vagna ok veðra} ‘the land of sun and moon, and the heavenly bodies, wagons and winds.’}} \hld\ ek em á \alst{v}it kominn &
\ind bręgði ęngi \alst{f}ǫstu hęiti \alst{f}ira.“\eva

\bvb\speakernoteb{[Quoth Allwise:]}%
“Allwise I am called; I live beneath the earth; \\
\ind I own under a stone my home. \\
The man of wagons \ken*{= Thunder} I am come to visit; \\
\ind let no man break a firm vow!”\evb\evg


\bvg\bva\mssnote{\Regius~19v/12--14}%
\speakernote{[Þórr kvað:]}%
\llap{„}Ek mun \alst{b}ręgda \hld\ því’t ek \alst{b}rúðar ȧ &
\ind \alst{f}lęst umb rǫ́ð sem \alst{f}aðir. &
vas’k⸗a ek \alst{h}ęima \hld\ þȧ’s þér \alst{h}ęitit vas &
\ind at sá ęinn es \alst{g}jǫf es með \alst{g}oðum.“\eva

\bvb\speakernoteb{[Quoth Thunder:]}%
“\emph{I} will break it, for about the bride \\
\ind I have the greatest say, as her father. \\
I was not at home when it was vowed to thee, \\
\ind but he [I] alone is the giver among the gods.”\evb\evg


\bvg\bva\mssnote{\Regius~19v/14--16}%
\speakernote{[All·víss kvað:]}%
\llap{„}Hvat ’s þat \alst{r}ekka \hld\ es ï \alst{r}ǫ́ðum tęlsk &
\ind \alst{f}ljóðs ins \alst{f}agr-glóa; &
\alst{f}jarra-\alst{f}lęina \hld\ þik munu \alst{f}áir kunna; &
\ind hvęrr hęfir þik \alst{b}augum \alst{b}orit?“\eva

\bvb\speakernoteb{[Quoth Allwise:]}%
“What sort of champion is this who claims to have a say \\
\ind about the fair-glowing girl? \\
O foreign tramp, few men will know thee; \\
\ind who has borne bighs to thee?”\evb\evg


\bvg\bva\mssnote{\Regius~19v/16--18}%
\speakernote{[Þórr kvað:]}%
\llap{„}\alst{V}ing-Þȯrr ek hęiti \hld\ ek hęfi \alst{v}íða ratat &
\ind \alst{s}onr em’k \alst{S}íð·grana; &
at \alst{ȯ}·sátt mïnni \hld\ skalt þat it \alst{u}nga man hafa &
\ind ok þat \alst{g}jaf-orð \alst{g}eta.“\eva

\bvb\speakernoteb{[Quoth Thunder:]}%
“\inx[P]{Wing-Thunder} I am called; I have widely roamed; \\
\ind I am the son of Sidegrane. \\
Against my assent shalt thou have this young girl, \\
\ind and get that gift-word!”\evb\evg


\bvg\bva\mssnote{\Regius~19v/18--20}%
\speakernote{[All·víss kvað:]}%
\llap{„}\alst{S}áttir þïnar \hld\ es ek vil \alst{s}nemma hafa &
\ind ok þat \alst{g}jaf-orð \alst{g}eta. &
\alst{ęi}ga vilja \hld\ heldr an \alst{á}n vera &
\ind \edtrans{þat it \alst{m}jall-hvíta \alst{m}an}{this snow-white girl}{\Bfootnote{For whiteness as a symbol of nobility cf.~Introduction to \textlink{Rigsthula}.}}.“\eva

\bvb\speakernoteb{[Quoth Allwise:]}%
“Thy assents I wish to have soon, \\
\ind and get that gift-word, \\
I would rather have than be without \\
\ind this snow-white girl.”\evb\evg


\bvg\bva\mssnote{\Regius~19v/20--21}%
\speakernote{[Þórr kvað:]}%
\llap{„}\alst{M}ęyjar ǫ̇stum \hld\ \alst{m}un⸗a þér verða &
\ind \alst{v}ísi gęstr of \alst{v}arit, &
ef þú ór \alst{h}ęimi kant \hld\ \alst{h}vęrjum at sęgja &
\ind allt þat’s ek \alst{v}il \alst{v}ita.\eva

\bvb\speakernoteb{[Quoth Thunder:]}%
“The maiden’s love will not be, \\
\ind O wise guest, denied thee, \\
if thou from every home canst tell \\
\ind all I wish to know:\evb\evg


\bvg\bva\mssnote{\Regius~19v/21--23}%
Sęg-ðu mér þat \alst{A}l·víss \hld\ \edtrans{\alst{ǫ}ll of rǫk fira}{all rakes of men}{\Bfootnote{To be understood as ‘all lore’.}} &
\ind \alst{v}ǫr⸗umk dvergr at \alst{v}itir, &
hvé sú \alst{jǫ}rð hęitir \hld\ es liggr fyr \alst{a}lda sonum &
\ind \alst{h}ęimi \alst{h}vęrjum ï.“\eva

\bvb Tell me this, Allwise---of all rakes of men, \\
\ind I think, dwarf, thou dost know: \\
what the earth is called which lies before the sons of men \\
\ind in every home.”\evb\evg


\bvg\bva\mssnote{\Regius~19v/23--25}%
\speakernote{[All·víss kvað:]}%
\llap{„}\alst{Jǫ}rð hęitir með mǫnnum \hld\ ęn með \alst{ǫ}lfum fold, &
\ind kalla \alst{v}ega \alst{v}anir. &
\alst{ï}-grǿn \alst{jǫ}tnar \hld\ \alst{a}lfar gróandi &
\ind kalla \alst{au}r \alst{u}pp-ręgin.“\eva

\bvb\speakernoteb{[Quoth Allwise:]}%
“Earth it is called among men but among Elves ‘fold’; \\
\ind call it ‘ways’ the Wanes, \\
‘evergreen’ Ettins, Elves ‘growing’; \\
\ind call it ‘mud’ the Up-reins.”\evb\evg


\bvg\bva\mssnote{\Regius~19v/25}%
\speakernote{[Þórr kvað:]}%
\llap{„}Sęg-ðu mér þat \alst{A}l·víss \hld\ \alst{ǫ}ll of rǫk fira &
\ind \alst{v}ǫr⸗umk dvergr at \alst{v}itir; &
hvé sá himinn hęitir \hld\ \edtrans{erakendi}{...}{\Bfootnote{A string too corrupt to restore without excessive conjecture; it at least appears to contain the relative pronoun \emph{er} ‘which’, younger form of \emph{es} and the adjective \emph{kęnndr} ‘known’.  Based on the first line, the alliteration must have fallen on \emph{h-}, and the root that first suggests itself is \emph{hę́ð} ‘height’.  A possible restoration is then \emph{es ȧ hę́ð es kęnndr} ‘which is known on high’.}} &
\ind \alst{h}ęimi \alst{h}vęrjum ï.“\eva

\bvb\speakernoteb{[Quoth Thunder:]}%
“Tell me this, Allwise---of all rakes of men, \\
\ind I think, dwarf, thou dost know: \\
what the heaven is called ... \\
\ind in every home.”\evb\evg


\bvg\bva\mssnote{\Regius~19v/26--27}%
\speakernote{[All·víss kvað:]}%
\llap{„}\alst{H}iminn hęitir með mǫnnum \hld\ ęn \alst{H}lýrnir með goðum &
\ind kalla \alst{V}ind-ófni \alst{v}anir; &
\alst{u}pp-hęim \alst{jǫ}tnar \hld\ \alst{a}lfar fagra-rę́fr &
\ind \alst{d}vergar \alst{d}rjúpan sal.“\eva

\bvb\speakernoteb{[Quoth Allwise:]}%
“Heaven it is called among men but ‘Lerner’ among the Gods; \\
\ind ‘Wind-ovner’ call it the Wanes; \\
‘upham’ Ettins, Elves ‘fair roof’, \\
\ind Dwarfs ‘dripping hall’.”\evb\evg


\bvg\bva\mssnote{\Regius~19v/27--28}%
\speakernote{[Þórr kvað:]}%
\llap{„}Sęg-ðu mér þat \alst{A}l·víss \hld\ \alst{ǫ}ll of rǫk fira &
\ind \alst{v}ǫr⸗umk dvergr at \alst{v}itir; &
hvęrsu \alst{m}ȧni hęitir \hld\ sá’s \alst{m}ęnn sjá &
\ind \alst{h}ęimi \alst{h}vęrjum ï.“\eva

\bvb\speakernoteb{[Quoth Thunder:]}%
“Tell me this, Allwise---of all rakes of men, \\
\ind I think, dwarf, thou dost know: \\
how the moon is called which men do see \\
\ind in every home.”\evb\evg


\bvg\bva\mssnote{\Regius~19v/28--30}%
\speakernote{[All·víss kvað:]}%
\llap{„}\alst{M}ȧni hęitir með \alst{m}ǫnnum \hld\ ęn \alst{M}ylinn með goðum, &
\ind kalla \alst{h}verfanda \alst{h}vél \alst{h}ęlju ï; &
\alst{sk}yndi jǫtnar \hld\ ęn \alst{sk}in dvergar &
\ind kalla \alst{a}lfar \edtrans{\alst{á}r-tala}{year-tallier}{\Bfootnote{The moon was important in the Germanic calendar (witness \emph{month}, a “moon-th”).  Cf.~\textlink{Voluspa}[6] and \textlink{Vafthrudnismal}[23], 25.}}.“\eva

\bvb\speakernoteb{[Quoth Allwise:]}%
“Moon it is called among men, but ‘Milen’ among the Gods, \\
\ind they call it ‘turning wheel’ in Hell, \\
‘haster’ Ettins and ‘shine’ Dwarfs; \\
\ind Elves call it ‘year-tallier’.”\evb\evg


\bvg\bva\mssnote{\Regius~19v/30}%
\speakernote{[Þórr kvað:]}%
\llap{„}Sęg-ðu mér þat \alst{A}l·víss \hld\ \alst{ǫ}ll of rǫk fira &
\ind \alst{v}ǫr⸗umk dvergr at \alst{v}itir; &
hvé sú \alst{s}ól hęitir \hld\ es \alst{s}já alda synir. &
\ind \alst{h}ęimi \alst{h}vęrjum ï.“\eva

\bvb\speakernoteb{[Quoth Thunder:]}%
“Tell me this, Allwise---of all rakes of men, \\
\ind I think, dwarf, thou dost know: \\
what the sun is called, which the sons of men see, \\
\ind in every home.”\evb\evg


\bvg\bva\mssnote{\Regius~19v/31--32}%
\speakernote{[All·víss kvað:]}%
\llap{„}\edtext{\alst{S}ól hęitir með mǫnnum \hld\ ęn \alst{s}unna}{\lemma{Sól \dots\ sunna ‘Sun \dots\ Swail’}\Bfootnote{A problem arises for the translation since \emph{sól} is the normal word for ‘sun’ in ON, whereas the direct cognate to English “sun” \emph{sunna} is a poetic synonym.
\emph{sól} and \emph{sunna} are in fact doublets, ultimately deriving from different declensions of the IE heteroclitic noun nom.~*séh₂u-el, gen.~*sh₂-en- \char`~\ *sh₂u-n-ós \parencite[463--464]{Kroonen2013}.  To retain this symmetry in the translation ON \emph{sunna} has been rendered as ‘swail’ (< OE \emph{sweġel}), a cognate of \emph{sól} that preserves the old nominative-accusative \emph{*-l} of the PIE noun.}} með goðum &
\ind kalla \alst{d}vergar \alst{D}valins lęika; &
\alst{Ęy}-glói \alst{jǫ}tnar \hld\ \alst{a}lfar fagra-hvél &
\ind \alst{a}l-skír \alst{ȧ}sa synir.“\eva

\bvb\speakernoteb{[Quoth Allwise:]}%
“Sun it is called among men, but ‘Swail’ among the Gods, \\
\ind Dwarfs call it ‘Dwollen’s playmate’; \\
‘ever-glower’ Ettins, Elves ‘fair wheel’, \\
\ind ‘all-bright’ the sons of the Eese.”\evb\evg


\bvg\bva\mssnote{\Regius~19v/32--33}%
\speakernote{[Þórr kvað:]}%
\llap{„}Sęg-ðu mér þat \alst{A}l·víss \hld\ \alst{ǫ}ll of rǫk fira &
\ind \alst{v}ǫr⸗umk dvergr at \alst{v}itir; &
hvé þau \alst{sk}ý hęita \hld\ es \alst{sk}úrum blanda⸗sk &
\ind \alst{h}ęimi \alst{h}vęrjum ï.“\eva

\bvb\speakernoteb{[Quoth Thunder:]}%
“Tell me this, Allwise---of all rakes of men, \\
\ind I think, dwarf, thou dost know: \\
what the clouds are called which are mixed with showers \\
\ind in every home.”\evb\evg


\bvg\bva\mssnote{\Regius~19v/33--20r/1}%
\speakernote{[All·víss kvað:]}%
\llap{„}\alst{Sk}ý hęita með mǫnnum, \hld\ ęn \alst{sk}úr-vǫ́n með goðum; &
\ind kalla \alst{v}ind-flot \alst{v}anir; &
\alst{ú}r-vǫ́n \alst{jǫ}tnar, \hld\ \alst{a}lfar veðr-męgin; &
\ind kalla ï hęlju \alst{h}jalm \alst{h}uliðs.“\eva

\bvb\speakernoteb{[Quoth Allwise:]}%
“Clouds they are called among men, but ‘shower-hope’ among the Gods; \\
\ind ‘wind-fat’ the Wanes call them; \\
‘drizzle-hope’ the Ettins, Elves ‘weather-strength’; \\
\ind in Hell they call them ‘helmet of the hidden’.”\evb\evg


\bvg\bva\mssnote{\Regius~20r/1}%
\speakernote{[Þórr kvað:]}%
\llap{„}Sęg-ðu mér þat \alst{A}l·víss \hld\ \alst{ǫ}ll of rǫk fira &
\ind \alst{v}ǫr⸗umk dvergr at \alst{v}itir; &
hvé sá \alst{v}indr hęitir \hld\ es \alst{v}íðast fęrr &
\ind \alst{h}ęimi \alst{h}vęrjum ï.“\eva

\bvb\speakernoteb{[Quoth Thunder:]}%
“Tell me this, Allwise---of all rakes of men, \\
\ind I think, dwarf, thou dost know: \\
what the wind is called where travels widest \\
\ind in every home.”\evb\evg


\bvg\bva\mssnote{\Regius~20r/2--3 \\ TODO}%
\Ballnote{This st.~is cited in \Skaldskaparmal\ 74.}%
\speakernote{[All·víss kvað:]}%
\llap{„}\alst{V}indr hęitir með mǫnnum, \hld\ ęn \alst{v}ǫ́fuðr með goðum; &
\ind kalla \alst{g}nęggjuð ginn-ręgin. &
\alst{ǿ}pi \alst{jǫ}tnar \hld\ \alst{a}lfar dyn-fara &
\ind kalla ï \alst{h}ęlju \alst{h}viðuð.“\eva

\bvb\speakernoteb{[Quoth Allwise:]}%
“Wind it is called among men but ‘waver’ among the Gods, \\
\ind ‘neigher’ call it the Yin-Reins; \\
‘weeper’ Ettins, Elves ‘din-farer’; \\
\ind in Hell they call it ‘stormer’.”\evb\evg


\bvg\bva\mssnote{\Regius~20r/3--4}%
\speakernote{[Þórr kvað:]}%
\llap{„}Sęg-ðu mér þat \alst{A}l·víss \hld\ \alst{ǫ}ll of rǫk fira &
\ind \alst{v}ǫr⸗umk dvergr at \alst{v}itir; &
hvé þat \alst{l}ogn hęitir \hld\ es \alst{l}iggja skal &
\ind \alst{h}ęimi \alst{h}vęrjum ï.“\eva

\bvb\speakernoteb{[Quoth Thunder:]}%
“Tell me this, Allwise---of all rakes of men, \\
\ind I think, dwarf, thou dost know: \\
what the calm is called, which shall lie \\
\ind in every home.”\evb\evg


\bvg\bva\mssnote{\Regius~20r/4--5}%
\speakernote{[All·víss kvað:]}%
\llap{„}\alst{L}ogn hęitir með mǫnnum, \hld\ ęn \alst{l}ę́gi með goðum, &
\ind \edtrans{kalla \alst{v}inds flot \alst{v}anir}{‘wind’s fat’ the Wanes call it}{\Bfootnote{Most likely corruption from st.~18/2 above.}}; &
\alst{o}f-hlý \alst{jǫ}tnar \hld\ \alst{a}lfar dag-sefa, &
\ind kalla \alst{d}vergar \alst{d}ags veru.“\eva

\bvb\speakernoteb{[Quoth Allwise:]}%
“Calm it is called among men and ‘lowering’ among the Gods, \\
\ind ‘wind’s fat’ the Wanes call it; \\
‘great lee’ Ettins, Elves ‘day-sleep’, \\
\ind call it Dwarfs ‘day’s rest’.”\evb\evg


\bvg\bva\mssnote{\Regius~20r/5}%
\speakernote{[Þórr kvað:]}%
\llap{„}Sęg-ðu mér þat \alst{A}l·víss \hld\ \alst{ǫ}ll of rǫk fira &
\ind \alst{v}ǫr⸗umk dvergr at \alst{v}itir; &
hvé sá \alst{m}arr hęitir \hld\ es \alst{m}ęnn róa &
\ind \alst{h}ęimi \alst{h}vęrjum ï.“\eva

\bvb\speakernoteb{[Quoth Thunder:]}%
“Tell me this, Allwise---of all rakes of men, \\
\ind I think, dwarf, thou dost know: \\
what the ocean is called where men do row \\
\ind in every home.”\evb\evg


\bvg\bva\mssnote{\Regius~20r/6--7}%
\speakernote{[All·víss kvað:]}%
\llap{„}\alst{S}ę́r hęitir með mǫnnum, \hld\ ęn \alst{s}ï-lę́gja með goðum, &
\ind kalla \alst{v}ág \alst{v}anir; &
\alst{á}l-hęim \alst{jǫ}tnar, \hld\ \alst{a}lfar laga-staf, &
\ind kalla \alst{d}vergar \alst{d}júpan mar.“\eva

\bvb\speakernoteb{[Quoth Allwise:]}%
“Sea it is called among men but ‘ever-low’ among the Gods; \\
\ind ‘wave’ the Wanes call it; \\
‘eelhome’ Ettins, Elves ‘staff of waters’; \\
\ind Dwarfs call it ‘deep ocean’.”\evb\evg


\bvg\bva\mssnote{\Regius~20r/7--8}%
\speakernote{[Þórr kvað:]}%
\llap{„}Sęg-ðu mér þat \alst{A}l·víss \hld\ \alst{ǫ}ll of rǫk fira &
\ind \alst{v}ǫr⸗umk dvergr at \alst{v}itir; &
hvé sá \alst{ę}ldr hęitir \hld\ es brenn fyr \alst{a}lda sonum &
\ind \alst{h}ęimi \alst{h}vęrjum ï.“\eva

\bvb\speakernoteb{[Quoth Thunder:]}%
“Tell me this, Allwise---of all rakes of men, \\
\ind I think, dwarf, thou dost know: \\
what the fire is called, which burns for the sons of men, \\
\ind in every home.”\evb\evg


\bvg\bva\mssnote{\Regius~20r/8--9}%
\speakernote{[All·víss kvað:]}%
\llap{„}\alst{Ę}ldr hęitir með mǫnnum \hld\ ęn með \alst{ǫ̇}sum funi &
\ind \edtrans{kalla \alst{v}ág \alst{v}anir}{‘wave’ the Wanes call it}{\Bfootnote{Certainly an error based on st.~24/2 above; fire is not otherwise known as ‘wave’.}}; &
\alst{f}rekan jǫtnar \hld\ ęn \alst{f}or-bręnni d*vergar &
\ind kalla ï \alst{h}ęlju \alst{h}rǫðuð.“\eva

\bvb\speakernoteb{[Quoth Allwise:]}%
“Fire it is called among men but among the Eese ‘flame’, \\
\ind ‘wave’ the Wanes call it; \\
‘the greedy’ Ettins, but ‘burner’ Dwarfs; \\
\ind in Hell they call it ‘hurrier’.”\evb\evg


\bvg\bva\mssnote{\Regius~20r/9--10}%
\speakernote{[Þórr kvað:]}%
\llap{„}Sęg-ðu mér þat \alst{A}l·víss \hld\ \alst{ǫ}ll of rǫk fira &
\ind \alst{v}ǫr⸗umk dvergr at \alst{v}itir; &
hvé \alst{v}iðr hęitir \hld\ es \alst{v}ęx fyr alda sonum &
\ind \alst{h}ęimi \alst{h}vęrjum ï.“\eva

\bvb\speakernoteb{[Quoth Thunder:]}%
“Tell me this, Allwise---of all rakes of men, \\
\ind I think, dwarf, thou dost know: \\
what the wood is called, which grows for the sons of men, \\
\ind in every home.”\evb\evg


\bvg\bva\mssnote{\Regius~20r/10--11}%
\speakernote{[All·víss kvað:]}%
\llap{„}\alst{V}iðr hęitir með mǫnnum. \hld\ ęn \edtext{\alst{v}allar fa\emph{x}}{\Afootnote{emend.; \emph{vallar far} \Regius.}} með goðum &
\ind kalla \alst{h}líð-þang \alst{h}alir; &
\alst{ę}ldi \alst{jǫ}tnar \hld\ \alst{a}lfar fagr-lima &
\ind kalla \alst{v}ǫnd \alst{v}anir.“\eva

\bvb\speakernoteb{[Quoth Allwise:]}%
“Wood it is called among men but ‘mane of the plain’ among the Gods; \\
\ind ‘slope-kelp’ heroes call it; \\
‘firewood’ Ettins, Elves ‘fair-limb’; \\
\ind ‘wand’ the Wanes call it.”\evb\evg


\bvg\bva\mssnote{\Regius~20r/11--12}%
\speakernote{[Þórr kvað:]}%
\llap{„}Sęg-ðu mér þat \alst{A}l·víss \hld\ \alst{ǫ}ll of rǫk fira &
\ind \alst{v}ǫr⸗umk dvergr at \alst{v}itir; &
hvé sú \alst{N}ǫ́tt hęitir \hld\ \edtrans{in \alst{N}ǫrvi kęnda}{born to Narrow}{\Bfootnote{See \textlink{Vafthrudnismal}[25]/2 n.}} &
\ind \alst{h}ęimi \alst{h}vęrjum ï.“\eva

\bvb\speakernoteb{[Quoth Thunder:]}%
“Tell me this, Allwise---of all rakes of men, \\
\ind I think, dwarf, thou dost know: \\
what the Night is called, born to Narrow, \\
\ind in every home.”\evb\evg


\bvg\bva\mssnote{\Regius~20r/12--13 \\ TODO}%
\Ballnote{This st.~is cited in \Skaldskaparmal\ 78.}%
\speakernote{[All·víss kvað:]}%
\llap{„}\alst{N}ǫ́tt hęitir með mǫnnum \hld\ ęn \alst{n}jól með goðum, &
\ind kalla \alst{g}rímu \alst{g}inn-ręgin; &
\alst{ȯ}·ljós \alst{jǫ}tnar \hld\ \alst{a}lfar svefn-gaman &
\ind kalla \alst{d}vergar \alst{d}raum-Njǫrun.“\eva

\bvb\speakernoteb{[Quoth Allwise:]}%
“Night it is called among men but ‘nivel’ among the Gods; \\
\ind ‘mask’ the Yin-Reins call it, \\
‘un-light’ Ettins, Elves ‘sleep-joy’; \\
\ind Dwarfs call it ‘dream-Narn’.”\evb\evg


\bvg\bva\mssnote{\Regius~20r/13--14}%
\speakernote{[Þórr kvað:]}%
\llap{„}Sęg-ðu mér þat \alst{A}l·víss \hld\ \alst{ǫ}ll of rǫk fira &
\ind \alst{v}ǫr⸗umk dvergr at \alst{v}itir; &
hvé þat \alst{s}ǫ́ð hęitir \hld\ es \alst{s}áa alda synir &
\ind \alst{h}ęimi \alst{h}vęrjum ï.“\eva

\bvb\speakernoteb{[Quoth Thunder:]}%
“Tell me this, Allwise---of all rakes of men, \\
\ind I think, dwarf, thou dost know: \\
what the seed is called, which the sons of men sow \\
\ind in every home.”\evb\evg


\bvg\bva\mssnote{\Regius~20r/14--15}%
\speakernote{[All·víss kvað:]}%
\llap{„}\alst{B}ygg hęitir með mǫnnum \hld\ ęn \alst{b}arr með goðum &
\ind kalla \alst{v}ǫxt \alst{v}anir. &
\alst{ę́}ti \alst{jǫ}tnar \hld\ \alst{a}lfar laga-staf &
\ind kalla ï \alst{h}ęlju \alst{h}nipinn.“\eva

\bvb\speakernoteb{[Quoth Allwise:]}%
“Barley it is called among men but ‘leaf’ among the Gods; \\
\ind ‘growth’ the Wanes call it; \\
‘eating’ Ettins, Elves ‘staff of waters’; \\
\ind in Hell they call it ‘wilting’.”\evb\evg


\bvg\bva\mssnote{\Regius~20r/15--16}%
\speakernote{[Þórr kvað:]}%
\llap{„}Sęg-ðu mér þat \alst{A}l·víss \hld\ \alst{ǫ}ll of rǫk fira &
\ind \alst{v}ǫr⸗umk dvergr at \alst{v}itir; &
hvé þat \alst{ǫ}l hęitir \hld\ es drekka \alst{a}lda synir &
\ind \alst{h}ęimi \alst{h}vęrjum ï.“\eva

\bvb\speakernoteb{[Quoth Thunder:]}%
“Tell me this, Allwise---of all rakes of men, \\
\ind I think, dwarf, thou dost know: \\
what the ale is called, which the sons of men drink, \\
\ind in every home.”\evb\evg


\bvg\bva\mssnote{\Regius~20r/16--17}%
\speakernote{[All·víss kvað:]}%
\llap{„}\alst{Ǫ}l hęitir með mǫnnum \hld\ ęn með \alst{ǫ̇}sum bjórr; &
\ind kalla \alst{v}ęig \alst{v}anir; &
\alst{h}ręina-lǫg jǫtnar \hld\ ęn ï \alst{h}ęlju mjǫð; &
\ind kalla \alst{s}umbl \alst{S}uttungs \alst{s}ynir.“\eva

\bvb\speakernoteb{[Quoth Allwise:]}%
“Ale it is called among men but among the Eese ‘beer’; \\
\ind ‘draughts’ the Wanes call it; \\
‘pure water’ the Ettins but in Hell ‘mead’; \\
\ind ‘simble’ Sutting’s sons call it.”\evb\evg


\bvg\bva\mssnote{\Regius~20r/18--19}%
\speakernote{[Þórr kvað:]}%
\llap{„}İ \alst{ęi}nu brjósti \hld\ ek sá’k \alst{a}ldri⸗gi &
\ind \alst{f}lęiri \alst{f}orna stafi; &
miklum \alst{t}ǫ́lum \hld\ ek kveð \alst{t}ę́ldan þik: &
\ind uppi est \alst{d}vergr of \alst{d}agaðr; &
\ind nú skïnn \alst{s}ól ï \alst{s}ali.“\eva

\bvb\speakernoteb{[Quoth Thunder:]}%
“In a single breast I never saw \\
\ind more ancient staves--- \\
with mighty tricks I call thee tricked: \\
\ind thou art, dwarf, dayed up; \\
\ind now shines the sun into the halls!”\evb\evg

\sectionline
